\paragraph{Autonomous collision clocks and their onset control.}
An even native charged-branch transfer, its Gauss reference and neutral
memory define a complete collision source with phase-insensitive pumping
and two-quantum loss. A three-state battery with energies $m,0,3m$ supplies
an energy-conserving dilation. Every event, inactive occupation and Fock
complement is retained. Exact positive-matrix bounds certify a coherence
frequency $m$ and a nonzero finite linewidth. For arrays of four, eight and
sixteen native branches the first-band linewidths are approximately
$0.0333783$, $0.0139327$ and $0.00602053$ in source energy units. The array
frequency is its constituent gap, not its total rest energy.

For a single charged branch the gap is the inherited particle rest energy.
Its complete equatorial records have the factor
$(1-\gamma\delta)^{k/2}\exp(-imk\delta)$. An exact counterfactual shows
that the first probe prepares this conditional phase independently of the
prior pumped occupation: with the pump removed, the oscillation remains
and its linewidth vanishes. Autonomous particle occupation onset is
therefore established, while autonomous internal phase onset is not
established by this instrument. This distinction concerns the actual
observer program; a vanishing ensemble mean by itself does not exclude
oscillatory shared conditional records.
